% Generalized Einsum with Row Variables:
% Shape Inference for Deep Learning in OCaml
%
% Workshop paper targeting OCaml Workshop 2026 (ICFP)

\documentclass[acmsmall,nonacm,review]{acmart}

\usepackage{listings}
\usepackage{xcolor}
\usepackage{booktabs}

% OCaml listing style
\lstdefinelanguage{ocaml}{
  morekeywords={let,in,fun,match,with,type,val,module,struct,sig,end,
    open,if,then,else,rec,and,or,not,true,false,of,Some,None,
    mutable,begin,end},
  morekeywords=[2]{Shape,Tensor,TDSL,NTDSL,PDSL,Ir},
  sensitive=true,
  morecomment=[s]{(*}{*)},
  morestring=[b]",
}

\lstset{
  language=ocaml,
  basicstyle=\small\ttfamily,
  keywordstyle=\bfseries\color{blue!70!black},
  keywordstyle=[2]\color{purple!70!black},
  commentstyle=\itshape\color{green!50!black},
  stringstyle=\color{red!70!black},
  columns=flexible,
  breaklines=true,
  showstringspaces=false,
  xleftmargin=1em,
  frame=single,
  framerule=0.4pt,
  numbers=none,
  escapeinside={(@}{@)},
  aboveskip=0.5\baselineskip,
  belowskip=0.5\baselineskip,
}

% Inline code
\newcommand{\code}[1]{\lstinline[basicstyle=\normalsize\ttfamily]{#1}}
\newcommand{\spec}[1]{\texttt{\small #1}}

\begin{document}

\title{Generalized Einsum with Row Variables:\\Shape Inference for Deep Learning in OCaml}

\author{Lukasz Stafiniak}
\affiliation{%
  \institution{ahrefs}
  \country{Poland}
}
\email{lukstafi@gmail.com}

\begin{abstract}
Tensor shape mismatches are among the most common bugs in deep learning code.
We present the shape inference system of OCANNL, an OCaml deep learning
framework that combines \emph{generalized einsum notation} with \emph{row
variables} and \emph{constraint-based inference}. Our einsum notation extends
the standard contraction syntax with strided iteration (for convolutions),
axis concatenation (for block tensors), and row variables (for flexible axis
counts). Shapes are organized into three axis kinds---batch, input, and
output---reflecting deep learning semantics. A multi-stage constraint solver
resolves dimension and row variables, producing informative error messages via
provenance tracking. We integrate this system into OCaml through PPX syntax
extensions, providing an ergonomic DSL that leverages OCaml's type and module
systems. We illustrate the approach with multi-head attention, convolutions,
and transformer blocks drawn from working OCANNL code.
\end{abstract}

\maketitle

%% ============================================================
\section{Introduction}
\label{sec:intro}

Deep learning frameworks must reconcile two competing demands: concise
notation for tensor operations, and early detection of shape errors.
Existing approaches each sacrifice one for the other:

\begin{itemize}
\item \textbf{Standard einsum} (NumPy, PyTorch) provides compact contraction
  notation but is limited to summation over shared indices. Convolutions,
  strided access, and concatenation require separate APIs. Shapes are checked
  at runtime.
\item \textbf{Named tensors} (PyTorch named tensors~\cite{chiang2021named},
  torchdim~\cite{ansel2024pytorch2}) attach semantic labels to axes, but
  labels are verbose in einsum-style notation and create inference ambiguity
  when multiple axes share the same semantic unit.
\item \textbf{Dependent types} (Dex~\cite{paszke2021getting},
  Futhark~\cite{henriksen2017futhark}) offer compile-time guarantees but
  demand heavy type annotations.
\item \textbf{einops}~\cite{rogozhnikov2022einops} provides readable
  reshaping syntax but operates at the level of shape transformations rather
  than computations.
\end{itemize}

OCANNL\footnote{\url{https://github.com/ahrefs/ocannl}} takes a different
approach. Its shape system combines three ideas:
\begin{enumerate}
\item A \emph{generalized einsum notation} that unifies contraction, strided
  iteration, convolution, and concatenation in a single syntax.
\item \emph{Row variables}, borrowed from row polymorphism in type theory,
  that allow specifications without committing to exact axis counts.
\item \emph{Constraint-based inference} with a multi-stage solver that
  resolves shapes at model construction time, before any data flows through
  the network.
\end{enumerate}

These are integrated into OCaml via PPX syntax extensions (\code{\%op} and
\code{\%cd}), yielding an embedded DSL that feels native while providing
shape safety. This paper describes the design, the inference algorithm, and
its application to standard deep learning patterns.

%% ============================================================
\section{The Shape System}
\label{sec:shapes}

\subsection{Three Axis Kinds}
\label{sec:axis-kinds}

Every tensor shape in OCANNL is partitioned into three rows of axes:
\[
  \underbrace{\text{batch}}_{\text{data samples}} \;\big|\;
  \underbrace{\text{input}}_{\text{consumed}} \to
  \underbrace{\text{output}}_{\text{produced}}
\]

This partition reflects deep learning semantics:
\begin{itemize}
\item \textbf{Batch axes} index independent data samples and are preserved
  across all operations.
\item \textbf{Output axes} represent produced values (leftmost in
  conventional matrix notation).
\item \textbf{Input axes} represent consumed values (rightmost), analogous to
  function arguments.
\end{itemize}

The input--output distinction gives tensor multiplication (\code{*}) the
semantics of function composition: for tensors with shapes $|i \to o_1$ and
$|i' \to o_2$, multiplication contracts $i$ against $o_2$, producing
$|i' \to o_1$---matching matrix multiplication when each row has a single
axis. This convention extends naturally to higher-order tensors without
requiring explicit einsum specs for common operations.

Each row is independently broadcastable: a dimension-1 axis broadcasts to
match any dimension, and a row with fewer axes is extended with dimension-1
axes. Unlike NumPy, which pads from the left, OCANNL broadcasts ``in the
middle''---preserving both specified leading and trailing axes of a row.

\subsection{Generalized Einsum Notation}
\label{sec:einsum}

OCANNL's einsum notation extends the standard contraction syntax in several
directions. The general form for binary operations is:

\begin{center}
\spec{rhs1\_spec ; rhs2\_spec => lhs\_spec}
\end{center}

\noindent where each spec describes the batch, input, and output axes of the
corresponding tensor using the \spec{batch | input -> output} convention.
Unary operations use \spec{rhs\_spec => lhs\_spec}. Within each spec, axes
are identified positionally by single-character or multi-character
labels.\footnote{Multi-character mode is triggered by the presence of a comma:
\spec{"abc"} denotes three axes \spec{a}, \spec{b}, \spec{c}, while
\spec{"abc,"} denotes a single axis named \spec{abc}.}

\paragraph{Standard contraction.}
Matrix multiplication is written as:
\begin{lstlisting}
a +* "m n; n p => m p" b
\end{lstlisting}
The \code{+*} operator performs addition-reduction with multiplication (the
standard einsum semantics). Labels appearing in both operand specs but not in
the result spec are contracted.

\paragraph{Strided iteration and convolutions.}
Affine index expressions in specs support strided access:
\begin{lstlisting}
x +* "... | stride*oh+kh, stride*ow+kw, ic;
      kh, kw, ic -> oc
      => ... | oh, ow, oc" ["kh"; "kw"] { kernel }
\end{lstlisting}
The expression \spec{stride*oh+kh} means the input index is
$\mathit{stride} \times \mathit{oh} + \mathit{kh}$, where \spec{oh} is the
output position and \spec{kh} is the kernel position. Adding a dilation
factor gives the standard convolution formula:
$\mathit{stride} \times o + \mathit{dilation} \times k$. The shape inference
system automatically derives the output dimensions from the input dimensions,
stride, dilation, and kernel size via the relationship:
\[
  d_{\text{in}} = \mathit{stride} \times (d_{\text{out}} - 1) +
  1 + (\mathit{d_{\text{kernel}}} - 1) \times \mathit{dilation}
\]

\paragraph{Concatenation.}
The \spec{\^{}} operator creates a concatenated axis:
\begin{lstlisting}
(* Concatenate two vectors *)
(a, b) ++^ "x; y => x^y"
(* Extract a prefix *)
x ++^ "a^b => a"
\end{lstlisting}
A concatenated axis \spec{a\^{}b} has dimension $d_a + d_b$ and iterates first
over the range of \spec{a}, then over the range of \spec{b}. This enables
block tensor construction and axis slicing within the same notation.

\paragraph{Dimension capture.}
Einsum specs can capture dimension variables for use in subsequent
computations:
\begin{lstlisting}
let scores =
  q +* "... s | h d; ... t | h d
        => ... s | t -> h" ["h"; "d"] k
in
Shape.set_dim h num_heads;
scores /. sqrt (dim d)
\end{lstlisting}
The list \code{["h"; "d"]} binds OCaml variables \code{h} and \code{d} to
the inferred dimensions of the corresponding axes. These can be used to set
dimension constraints (\code{Shape.set\_dim}) or to compute scaling factors
(\code{dim d} retrieves the integer dimension value).

\paragraph{Additional operators.}
Beyond the standard addition-multiply einsum (\code{+*}), OCANNL provides:
\texttt{@\^{}+} for tropical semiring operations (max-reduce with addition),
\code{++} for unary add-reduction, and \texttt{@\^{}\^{}} for unary
max-reduction. These share the same spec syntax, differing only in the
accumulation and combination operations.

\subsection{Row Variables}
\label{sec:row-vars}

Row variables are the key mechanism for writing shape-polymorphic
specifications. Borrowed from row polymorphism in type
theory~\cite{remy1993type}, they represent an unknown sequence of axes.

\begin{itemize}
\item \spec{...} is a context-dependent ellipsis that introduces an anonymous
  row variable for the enclosing axis kind.
\item \spec{..name..} (or \spec{..n..} in single-character mode) is a named
  row variable. Occurrences of the same name within a spec denote the same
  row variable and must unify to the same sequence of axes.
\end{itemize}

Row variables implement the ``principle of least commitment'': a
specification should constrain only the axes it cares about. For example,
the softmax spec:
\begin{lstlisting}
x ++ " ... | t -> ... => ... | 0 -> ..."
\end{lstlisting}
reduces over the output axis named \spec{t} while preserving all batch axes,
all other output axes, and all input axes---regardless of how many there are.
Without row variables, this would require separate implementations for each
possible axis count.

Row variables in OCANNL interact with broadcasting: a row variable can unify
with zero or more axes, and dim-1 axes introduced by broadcasting count as
legitimate axes. The inference system resolves row variables progressively:
first unifying what is determined by constraints, then closing remaining
variables to their least upper bounds or to ``no further axes'' at later
stages.

\subsection{Dimension Units}
\label{sec:units}

Axes have optional semantic annotations called \emph{dimension units}
(termed ``basis'' in the codebase). Unlike named tensors, where axes are
identified by their names, OCANNL axes are identified \emph{positionally}
and may optionally carry a unit that constrains matching: axes with different
units cannot be unified.

This design avoids two problems with named axes:
\begin{enumerate}
\item \textbf{Verbosity}: named axes require explicit labels on every axis of
  every tensor. Positional identification builds on mathematical tradition
  (linear algebra, Einstein notation) where position carries meaning.
\item \textbf{Inference ambiguity}: when multiple axes share the same
  semantic unit (e.g., two spatial dimensions both representing ``position''),
  named-axis systems cannot determine which name goes where. Positional
  identification with optional units avoids this entirely.
\end{enumerate}

For example, a unit \spec{"rgb"} attached to a dimension of size~3 denotes
that an axis represents color channels. Shape inference ensures that axes
unified with this one also have unit \spec{"rgb"} and size~3, providing
semantic safety. But the axis is still identified by its position in the
batch/input/output partition, not by its name.

%% ============================================================
\section{Constraint-Based Shape Inference}
\label{sec:inference}

\subsection{Representing Constraints}
\label{sec:constraints}

Shape inference maintains a global environment mapping dimension variables
and row variables to their current knowledge:

\begin{lstlisting}[language={},basicstyle=\small\ttfamily]
dim_var (@$\mapsto$@) Solved(dim) | Bounds(lub, constraints, ordering)
row_var (@$\mapsto$@) Solved(row) | Bounds(lub, constraints, ordering)
\end{lstlisting}

\noindent The constraint language includes:
\begin{itemize}
\item \textbf{Equations}: $d_1 = d_2$ (dimension equality), $r_1 = r_2$ (row
  equality).
\item \textbf{Inequalities}: $d_c \geq d_s$ (where $d_c$ is from a
  super-tensor and $d_s$ from a sub-tensor, reflecting broadcasting),
  $r_c \geq r_s$ (row inequality).
\item \textbf{Dimension constraints}: minimum dimension, total element count.
\item \textbf{Row constraints}: total element count (possibly involving
  strided dimension variables), exact axis lists.
\end{itemize}

For ground dimensions, $n \geq m$ means either $n = m$ or $m = 1$
(broadcasting). For ground rows, $q \geq r$ means $q$ has at least as many
axes as $r$, and for each corresponding pair of dimensions, $q$'s dimension
is $\geq$ $r$'s. This combines row polymorphism with nominal subtyping: the
subtyping relation states that dimension-1 (no unit) is smaller than all
dimensions.

Einsum operations generate \emph{equations} (not inequalities), so they do
not permit broadcasting---they require exact shape matching. Pointwise and
compose operations generate inequalities, allowing broadcasting.

\subsection{The Multi-Stage Solver}
\label{sec:stages}

Constraints are solved in stages, progressing from conservative (online,
as tensors are composed) to aggressive (on demand, when dimensions or
projections are needed):

\begin{enumerate}
\item \textbf{Stage~1 (online)}: Unification and constraint propagation as
  tensor expressions are constructed. This resolves what can be determined
  immediately.
\item \textbf{Stage~2}: Force precision-related coefficients (byte sizes of
  numeric types).
\item \textbf{Stage~3}: Close row variables in \emph{terminal} (leaf) shapes
  to their least upper bounds (LUBs). Terminal shapes are closed first
  because their shapes can only be determined by downstream constraints (they
  can always be broadcast), unlike composite shapes which are forced ``from
  below'' by their components.
\item \textbf{Stage~4}: Close terminal dimension variables with multiple
  lower bounds to their LUBs. Resolve \code{Total\_elems} constraints where
  row variables have been eliminated.
\item \textbf{Stage~5}: Guess remaining terminal dimension variables to~1
  (the broadcasting identity), unless prevented by \code{Total\_elems}
  constraints. Resolve remaining row constraints heuristically.
\item \textbf{Stage~6}: Extend LUB processing to non-terminal shapes. Set
  remaining row variables to ``no further axes.''
\item \textbf{Stage~7}: Close all remaining dimension variables to their
  lower bounds or to~1.
\end{enumerate}

The terminal-first strategy (Stages 3--5 before 6--7) reflects an asymmetry:
leaf tensors acquire shape information only from downstream constraints and
broadcasting, so they must be resolved first. Composite tensors then get
their shapes forced by their already-resolved components.

\paragraph{Preventing premature guessing.}
A subtle issue arises with strided dimensions (convolutions): if a dimension
variable appears in the numerator of a \code{Total\_elems} constraint (e.g.,
$\text{total} = \text{coeff} \times d / \text{denom}$), guessing $d = 1$
prematurely can make the constraint unsatisfiable. The solver tracks these
dependencies via a \code{has\_uniq\_constr\_unless} field: a numerator
variable cannot be guessed to~1 unless at least one denominator variable is
also prevented from guessing (breaking potential cycles).

\subsection{Error Provenance}
\label{sec:provenance}

Each row carries a \emph{provenance} record tracking which einsum spec, tensor
expression, or operation produced it. When unification fails---e.g., two axes
have incompatible dimensions or units---the error message includes the
provenance chain, allowing the user to trace the conflict back to the
source-level einsum specification.

%% ============================================================
\section{OCaml Integration}
\label{sec:ocaml}

OCANNL integrates with OCaml via two PPX syntax extensions:
\begin{itemize}
\item \code{\%op}: for tensor expressions with automatic differentiation.
  Records like \code{\{ w \}} declare learnable parameters.
\item \code{\%cd}: for assignment computations (code generation without
  gradient tracking).
\end{itemize}

Both extensions transform OCaml expressions, reinterpreting arithmetic
operators as tensor operations and record syntax as tensor declarations. A
unit parameter \code{()} separates configuration arguments (processed as
plain OCaml) from tensor arguments (transformed by the extension):

\begin{lstlisting}
let%op mlp_layer ~label ~hid_dim () x =
  relu (({ w } * x) + { b = 0.; o = [hid_dim] })
\end{lstlisting}

The record \code{\{ w \}} declares a parameter with default initialization;
\code{\{ b = 0.; o = [hid\_dim] \}} declares a bias initialized to zero with
explicit output dimensions. The \code{[hid\_dim]} syntax specifies a
one-dimensional output of size \code{hid\_dim}. The \code{\%oc} anti-quotation
escapes back to plain OCaml when needed inside a transformed expression.

This design leverages OCaml's module system: the extensions require specific
DSL modules in scope (\code{TDSL} for \code{\%op}, \code{NTDSL} for
\code{\%cd}), which determine gradient tracking behavior. The modules are
brought into scope uniformly via \code{open Ocannl.Operation.DSL\_modules}.

%% ============================================================
\section{Examples}
\label{sec:examples}

\subsection{Multi-Head Attention}
\label{sec:attention}

The multi-head attention mechanism illustrates several features working
together. In OCANNL:

\begin{lstlisting}
let%op multi_head_attention ~label ~num_heads
    ~d_k ~d_v ?temperature
    ?(dropout_rate = 0.0) ()
    ~train_step ?mask x =
  let q = { w_q } * x in
  let k = { w_k } * x in
  let v = { w_v } * x in
  let scores =
    (q +* k " ... s | h d; ... t | h d
             => ... s | t -> h" ["h"; "d"])
    /. sqrt (dim d)
  in
  Shape.set_dim h num_heads;
  Shape.set_dim d d_k;
  Shape.set_dim e d_v;
  let attn_weights =
    softmax ~spec:" ... | t -> ..." ?temperature ()
      (match mask with
       | None -> scores
       | Some mask -> where mask scores !.(-1e9))
  in
  let attn_weights =
    dropout ~rate:dropout_rate ()
      ~train_step attn_weights
  in
  { w_o } *
    (attn_weights +* v
       " ... s | t -> h; ... t | h e
         => ... s | h e" ["e"])
\end{lstlisting}

Key observations:
\begin{itemize}
\item \textbf{Row variables} (\spec{...}): the batch dimensions are handled
  polymorphically---the same code works for any number of batch axes.
\item \textbf{Input--output split}: in the score computation,
  \spec{... s | t -> h} makes \spec{t} (time/key positions) an output axis
  and \spec{h} (heads) an input axis. This means the softmax reduces over
  \spec{t} while keeping \spec{h} independent---exactly the attention
  semantics.
\item \textbf{Dimension capture}: \code{["h"; "d"]} captures the head and
  key dimensions for scaling (\code{sqrt (dim d)}) and constraint setting
  (\code{Shape.set\_dim h num\_heads}).
\item \textbf{Implicit shapes}: the projection matrices \code{\{w\_q\}},
  \code{\{w\_k\}}, \code{\{w\_v\}}, \code{\{w\_o\}} have their shapes
  entirely inferred from context---no dimension annotations needed.
\end{itemize}

\subsection{Convolution}
\label{sec:conv}

A 2D convolution with configurable stride and kernel size:

\begin{lstlisting}
let%op conv2d ?(stride=1) ?(kernel_size=3) () x =
  Shape.set_dim kh kernel_size;
  Shape.set_dim kw kernel_size;
  x +* "... | stride*oh+kh, stride*ow+kw, ic;
        kh, kw, ic -> oc
        => ... | oh, ow, oc" ["kh"; "kw"] { kernel }
  + { bias = 0. }
\end{lstlisting}

The spec \spec{stride*oh+kh} expresses that input position =
$\mathit{stride} \times \mathit{oh} + \mathit{kh}$. The inference system
derives the output spatial dimensions from the input dimensions, stride, and
kernel size. Row variables (\spec{...}) handle arbitrary batch dimensions.

\subsection{Layer Normalization}
\label{sec:layernorm}

Layer normalization demonstrates dimension capture and reduction:

\begin{lstlisting}
let%op layer_norm ~label ?(epsilon = 1e-5) () x =
  let mean =
    x ++ " ... | ..d..  => ... | 0 " ["d"] in
  let centered = (x - mean) /. dim d in
  let variance =
    (centered *. centered)
      ++ " ... | ... => ... | 0 " in
  let std_dev = sqrt (variance + !.epsilon) in
  let normalized = centered /. std_dev in
  ({ gamma = 1. } *. normalized) + { beta = 0. }
\end{lstlisting}

The spec \spec{... | ..d.. => ... | 0} reduces \emph{all} output axes to a
scalar (0 denotes a removed axis), while the row variable \spec{..d..}
captures the total output dimensionality for normalization. The named row
variable means that however many output axes there are, they are all reduced
and the product of their sizes is accessible via \code{dim d}.

\subsection{Comparison with Other Frameworks}
\label{sec:comparison}

Table~\ref{tab:comparison} compares how a dot-product attention score
computation is expressed across frameworks.

\begin{table}[t]
\caption{Attention score computation across frameworks.}
\label{tab:comparison}
\small
\begin{tabular}{@{}lp{5.5cm}@{}}
\toprule
\textbf{Framework} & \textbf{Expression} \\
\midrule
NumPy &
\texttt{np.einsum("bhsd,bhtd->bhst",\linebreak q, k) / sqrt(d\_k)} \\[3pt]
PyTorch &
\texttt{torch.matmul(q, k.transpose(-2,-1))\linebreak / math.sqrt(d\_k)} \\[3pt]
einops &
(no native einsum; reshape +\linebreak\texttt{torch.matmul}) \\[3pt]
OCANNL &
\texttt{(q +* k "...s|hd; ...t|hd\linebreak => ...s|t->h" ["h";"d"])\linebreak /. sqrt (dim d)} \\
\bottomrule
\end{tabular}
\end{table}

Key differences:
\begin{itemize}
\item NumPy einsum requires the user to commit to exact axis counts
  (\spec{bhsd})---adding a batch dimension requires changing the spec.
  OCANNL's \spec{...} handles this polymorphically.
\item PyTorch uses index-based transposition, obscuring the operation's intent.
  OCANNL's spec makes the contraction pattern explicit.
\item The OCANNL version captures \code{d} for scaling and \code{h} for
  setting the number of heads, avoiding separate bookkeeping.
\item The input/output distinction (\spec{t -> h} in the result) carries
  semantic meaning: \spec{t} is an output axis (produced by the operation),
  while \spec{h} is an input axis (preserved for the subsequent
  value-weighting step).
\end{itemize}

%% ============================================================
\section{Related Work}
\label{sec:related}

\paragraph{Einsum extensions.}
The original einsum notation dates to NumPy~\cite{harris2020array} and has
been adopted by TensorFlow and PyTorch. \texttt{opt\_einsum}~\cite{smith2018opt}
optimizes contraction order but does not extend the notation.
einops~\cite{rogozhnikov2022einops} extends einsum with reshape and repeat
operations but does not support convolutions or constraint-based inference.
OCANNL's generalized einsum subsumes standard contraction and adds strided
iteration, convolution, and concatenation.

\paragraph{Named tensors and axis labels.}
PyTorch named tensors~\cite{chiang2021named} and
torchdim/DumPy~\cite{ansel2024pytorch2} associate names with axes.
Named axes improve readability but introduce inference challenges: when
axes share semantic units, the mapping from names to positions becomes
ambiguous. OCANNL's positional axes with optional units avoid this while
providing semantic constraints.

\paragraph{Dependent types for arrays.}
Dex~\cite{paszke2021getting} uses dependent types with index sets to
ensure shape safety. Futhark~\cite{henriksen2017futhark} uses size types.
Both provide strong static guarantees but require explicit type annotations
that OCANNL's inference makes unnecessary. OCANNL's approach is closer in
spirit to Hindley-Milner inference: shapes are inferred from usage rather
than declared.

\paragraph{Row polymorphism.}
Row polymorphism was introduced by R{\'e}my~\cite{remy1993type} for record
types. It has been applied to effect systems~\cite{leijen2014koka} and
extensible records. To our knowledge, OCANNL is the first system to apply
row polymorphism to tensor shape inference, using row variables to represent
unknown sequences of axes.

%% ============================================================
\section{Discussion and Future Work}
\label{sec:discussion}

\paragraph{Limitations.}
OCANNL's shape system is currently monomorphic: variables are interpreted
existentially, not universally. A tensor-producing function that is called
twice shares shape variables between calls, which can over-constrain shapes.
The workaround is the unit-parameter pattern (\code{let layer = make () in
...}), which creates fresh variables per call. True shape polymorphism would
require abstracting over fresh variables when forming tensor functions and
freshening them at application sites---essentially an abstract interpretation
mechanism. This is planned for a future version.

\paragraph{Shape schemes.}
A natural extension would be \emph{shape schemes}: universally quantified
shape specifications that are freshened at each use site, analogous to
type schemes in Hindley-Milner. This would eliminate the need for
unit-parameter workarounds and enable shape-polymorphic libraries.

\paragraph{Formal guarantees.}
The current system relies on testing rather than formal proofs for
correctness. Formalizing the constraint language and proving soundness of
the multi-stage solver would strengthen confidence, particularly around the
interaction of row variables with broadcasting and strided dimensions.

%% ============================================================
\section{Conclusion}
\label{sec:conclusion}

We have presented OCANNL's shape inference system, which combines generalized
einsum notation, row variables, and constraint-based inference to provide
shape safety for deep learning code. The three axis kinds
(batch/input/output) capture deep learning semantics naturally. Row variables
enable the ``principle of least commitment''---specifications constrain only
the axes that matter. The multi-stage solver resolves shapes before execution,
with provenance-tracked error messages. Integration via OCaml PPX extensions
provides an ergonomic surface syntax. Our examples demonstrate that this
approach handles standard deep learning patterns (attention, convolutions,
normalization) concisely and safely.

OCANNL is open source and available at
\url{https://github.com/ahrefs/ocannl}.

%% ============================================================
\bibliographystyle{ACM-Reference-Format}

\begin{thebibliography}{10}

\bibitem{chiang2021named}
David Chiang, Alexander M.\ Rush, and Sasha Rush.
\newblock Named tensor notation.
\newblock In \emph{NeurIPS}, 2021.

\bibitem{ansel2024pytorch2}
Jason Ansel et al.
\newblock {PyTorch 2}: Faster machine learning through dynamic {Python}
  bytecode transformation and graph compilation.
\newblock In \emph{ASPLOS}, 2024.

\bibitem{paszke2021getting}
Adam Paszke, Daniel~D.\ Johnson, David Duvenaud, Dimitrios Vytiniotis,
  Alexey Radul, Matthew~J.\ Johnson, Jonathan Ragan-Kelley, and Dougal
  Maclaurin.
\newblock Getting to the point: Index sets and parallelism-preserving
  autodiff for pointful array computation.
\newblock \emph{Proc. ACM Program. Lang.}, 5(ICFP), 2021.

\bibitem{henriksen2017futhark}
Troels Henriksen, Niels G.~W.\ Serup, Martin Elsman, Fritz Henglein, and
  Cosmin~E.\ Oancea.
\newblock Futhark: Purely functional {GPU}-programming with nested
  parallelism and in-place array updates.
\newblock In \emph{PLDI}, 2017.

\bibitem{rogozhnikov2022einops}
Alex Rogozhnikov.
\newblock Einops: Clear and reliable tensor manipulations with
  Einstein-like notation.
\newblock In \emph{ICLR}, 2022.

\bibitem{remy1993type}
Didier R{\'e}my.
\newblock Type inference for records in a natural extension of {ML}.
\newblock In Carl~A.\ Gunter and John~C.\ Mitchell, editors, \emph{Theoretical
  Aspects of Object-Oriented Programming}, pages 67--95. MIT Press, 1993.

\bibitem{harris2020array}
Charles~R.\ Harris et al.
\newblock Array programming with {NumPy}.
\newblock \emph{Nature}, 585(7825):357--362, 2020.

\bibitem{smith2018opt}
Daniel G.~A.\ Smith and Johnnie Gray.
\newblock opt\_einsum: A {Python} package for optimizing contraction order
  for einsum-like expressions.
\newblock \emph{JOSS}, 3(26):753, 2018.

\bibitem{leijen2014koka}
Daan Leijen.
\newblock Koka: Programming with row polymorphic effect types.
\newblock \emph{MSFP}, 153:100--126, 2014.

\end{thebibliography}

\end{document}
